\section*{Capítulo \ref{ch:b2:euler-bernoulli} --- Fluidos perfeitos: Euler e Bernoulli}

\begin{solution}{exo:b2:euler-bernoulli:1}
$v_2 = 1.5 \times 4 = \qty{6.0}{m/s}$, $P_2 = 3.0 \times 10^5 - 500(36 - 2.25) =
\qty{2.83}{bar}$. A \qty{2}{cm}: $v = \qty{24}{m/s}$, $P = 3.0 \times 10^5 - 500
(576 - 2.25) = \qty{0.13}{bar}$ --- perto da pressão de vapor. Um resultado
negativo significa que a vazão suposta é impossível: a água cavita
e o escoamento é estrangulado.
\end{solution}

\begin{solution}{exo:b2:euler-bernoulli:2}
$v = \sqrt{2 \times 15000/0.74} = \qty{201}{m/s}$; indicada, $\sqrt{2 \times 15000/
1.22} = \qty{157}{m/s}$. A velocidade indicada mede a \omterm{thm:b2:euler-bernoulli:bernoulli}{pressão dinâmica},
que é o que a asa sente: a celeridade de estol é uma velocidade
\emph{indicada} fixa em qualquer altitude.
\end{solution}

\begin{solution}{exo:b2:euler-bernoulli:3}
Profundidade \qty{2.0}{m}: $v = \sqrt{2g \times 2} = \qty{6.3}{m/s}$; $D_V = 0.6 \times 10^{-4}
\times 6.3 = \qty{0.38}{L/s}$; tempo de queda $\sqrt{2/g} = \qty{0.45}{s}$, alcance
\qty{2.8}{m}. O alcance $= 2\sqrt{h(H - h)}$ é simétrico em $h \leftrightarrow H -
h$: um furo à profundidade \qty{1.0}{m} (\qty{2}{m} acima do piso) atinge o mesmo
ponto.
\end{solution}

\begin{solution}{exo:b2:euler-bernoulli:4}
$\Delta P = \tfrac12 \times 1.2 \times 900 = \qty{540}{Pa}$, com a pressão externa menor:
força resultante de \qty{54}{kN} \emph{para cima}; o telhado pesa \qty{49}{kN}: ele é
arrancado. (Telhados reais rompem primeiro nas bordas, onde o escoamento se descola.)
\end{solution}

\begin{solution}{exo:b2:euler-bernoulli:5}
(a) $\tan\theta = 3/9.81$, $\theta = \qty{17}{\degree}$; $2 \times 0.306 = \qty{0.61}{m}$
(mais alto atrás). (b) $\Delta P = \rho aL = \qty{6.0}{kPa}$. (c) $\tan\theta =
0.61$, diferença de \qty{1.22}{m}: a frente sobe \qty{0.61}{m} acima da média de
\qty{0.40}{m}: \qty{1.01}{m} $>$ \qty{0.8}{m}, e a água bate na tampa. (d) $a = v^2/R =
\qty{5.6}{m/s^2}$: $\theta = \qty{30}{\degree}$, mais alto do lado de fora da curva.
\end{solution}

\begin{solution}{exo:b2:euler-bernoulli:6}
(a) Paraboloide $z = z_0 + \Omega^2r^2/2g$, $\Omega = 4\pi = \qty{12.6}{rad/s}$. (b)
$\Omega^2R^2/2g = 158 \times 0.0225/19.6 = \qty{0.18}{m}$. (c) A altura média de um
paraboloide sobre o disco fica no meio: centro $0.20 - 0.09 = \qty{0.11}{m}$, borda
\qty{0.29}{m}. (d) O centro seca quando $\Omega^2R^2/4g = h_0$: $\Omega = \sqrt{4gh_0}/R =
\qty{18.7}{rad/s} = 3.0$ voltas por segundo. (e) Um paraboloide focaliza
raios paralelos num só ponto, em $f = g/2\Omega^2$; telescópios de espelho líquido
giram mercúrio.
\end{solution}

\begin{solution}{exo:b2:euler-bernoulli:7}
(a) $v = \sqrt{2g \times 3} = \qty{7.7}{m/s}$; $D_V = \pi \times 10^{-4} \times 7.7 =
\qty{2.4}{L/s}$. (b) $P = P_0 - \rho g(2) - \tfrac12\rho v^2 = 100 - 19.6 - 29.4 =
\qty{51}{kPa}$. (c) $P_0 - \rho gh_c - \tfrac12\rho v^2 \ge \qty{2.3}{kPa}$: $h_c \le
\qty{7.0}{m}$. (d) O ar é compressível e leve: não há coluna líquida
contínua para transmitir a pressão, e a água dos dois lados simplesmente fica onde está.
\end{solution}

\begin{solution}{exo:b2:euler-bernoulli:8}
(a) $\Delta P = mg/S = \qty{654}{Pa}$; $\tfrac12\rho v^2 = \qty{2.2}{kPa}$, $C_L = 0.30$.
(b) $v_{\text{sup}}^2 = 58^2 + 2 \times 654/1.22$: \qty{66.6}{m/s}. (c) $F_L' =
\qty{981}{N/m} = \rho v\Gamma$: $\Gamma = \qty{13.4}{m^2/s}$. (d) $v = \sqrt{2mg/
\rho SC_L} = \sqrt{19620/25.6} = \qty{28}{m/s}$.
\end{solution}

\begin{solution}{exo:b2:euler-bernoulli:9}
(a) $v_2 = \qty{2.5}{m/s}$; $\Delta P = \tfrac12 \times 1060 \times (6.25 - 0.25) =
\qty{3.2}{kPa}$, $P_2 = \qty{9.8}{kPa}$. (b) Ainda acima de \qty{1}{kPa}: não. (c) A
\qty{0.05}{cm^2}, $v = \qty{5}{m/s}$ e $\Delta P = \qty{13}{kPa}$: $P_2 \approx 0$, abaixo
da pressão do tecido --- a parede colapsa, o escoamento cessa, a
pressão se recupera e a artéria reabre: uma vibração (o sopro que o médico ouve).
\end{solution}

\begin{solution}{exo:b2:euler-bernoulli:10}
(a) $c_PT_0 = c_PT + v^2/2$ com $c_P = \gamma R/(\gamma - 1)M_{\text{mol}}$ e
$c^2 = \gamma RT/M_{\text{mol}}$: $v^2/2c_PT = \tfrac12(\gamma - 1)M^2$. (b) $M =
0.85$: $T_0 = 220 \times 1.145 = \qty{252}{K}$; $M = 2$: $220 \times 1.8 = \qty{396}{K}$;
$M = 25$: $220 \times 126 = \qty{28000}{K}$ --- sem sentido: o ar se dissocia
e se ioniza, e $c_P$ não é constante (as temperaturas de estagnação reais são
de vários milhares de kelvin). (c) $P_0/P = (1 + \tfrac{\gamma - 1}2M^2)^{\gamma/(\gamma
- 1)} \approx 1 + \tfrac{\gamma}2M^2 + \tfrac{\gamma}8M^4$, e $\gamma PM^2 = \rho v^2$:
$P_0 - P = \tfrac12\rho v^2(1 + M^2/4)$. Ler $v$ a partir de $\sqrt{2(P_0 - P)/\rho}$
a superestima em $\sqrt{1 + M^2/4} - 1 \approx 9\%$ em $M = 0.85$.
\end{solution}

\begin{solution}{exo:b2:euler-bernoulli:11}
(a) Da superfície ($v \approx 0$, $P_0$, altura $h$) até a saída ($v$, $P_0$,
altura $0$), com $\int\partial_tv\,\dd s = \ell\,\dd v/\dd t$. (b) $\dd v/(V^2 -
v^2) = \dd t/2\ell$, $V^2 = 2gh$: $\operatorname{artanh}(v/V) = Vt/2\ell$. (c) $V =
\qty{9.9}{m/s}$; $\tanh x = 0.9$ em $x = 1.47$: $t = 2\ell x/V = \qty{15}{s}$. (d)
$\dd v/\dd t \approx \qty{-20}{m/s^2}$: $\Delta P = \rho\ell \times 20 = \qty{9.9}{bar}$.
A água é ligeiramente compressível: um fechamento rápido lança uma onda de
pressão, $\Delta P = \rho c\Delta v \approx \qty{140}{bar}$.
\end{solution}

\begin{solution}{exo:b2:euler-bernoulli:12}
(a) Permanente, perfeito, incompressível e irrotacional: uma só constante
para todo o escoamento. (b) Na superfície $P = P_0$: $\tfrac12\rho v^2 + \rho gz =
0$, $z = -\Gamma^2/8\pi^2gr^2$. (c) Euler radial: $\partial_rP = \rho\Omega^2r$,
e hidrostática na vertical: a superfície satisfaz $z(r) - z(a) = \Omega^2(r^2 -
a^2)/2g$. (d) Banheira: $\Omega = \Gamma/2\pi a^2 = \qty{191}{rad/s}$; $z(a) = -\Omega^2a^2/
2g = \qty{-4.7}{cm}$, e o núcleo acrescenta outros \qty{4.7}{cm}: \qty{9.3}{cm} de profundidade.
Tornado: $\Delta P = \tfrac12\rho v_{\max}^2$ fora, mais o mesmo tanto dentro:
$\rho v_{\max}^2 = 1.2 \times 6400 = \qty{7.7}{kPa}$ abaixo da pressão ambiente no centro.
\end{solution}

\begin{solution}{pb:b2:euler-bernoulli:1}
\textbf{1.} $\rho g \times 150 = \qty{14.7}{bar}$ acima da atmosférica; $\tfrac12\rho gH_d^2
\times \qty{1}{m} = \qty{1.1e8}{N}$.

\textbf{2.} $\rho gH = \qty{19.6}{bar}$ efetivos, \qty{20.6}{bar} absolutos.

\textbf{3.} $v = 60/7.07 = \qty{8.5}{m/s}$; $\tfrac12\rho v^2 = \qty{0.36}{bar}$; $P =
P_0 + \rho gH - \tfrac12\rho v^2 = \qty{20.3}{bar}$ absolutos.

\textbf{4.} $v_j = \sqrt{2gH} = \qty{62.6}{m/s}$; $s = 60/62.6 = \qty{0.96}{m^2}$, $d =
\qty{1.1}{m}$.

\textbf{5.} $\tfrac12\rho D_Vv_j^2 = \tfrac12\rho D_V(2gH) = \rho gD_VH = \qty{118}{MW}$.

\textbf{6.} $m = \rho S\ell = \qty{2.8e6}{kg}$, $E_k = \tfrac12mv^2 = \qty{1.0e8}{J}$;
trânsito $\ell/v = \qty{47}{s}$.

\textbf{7.} $\rho gH = (P - P_0) + \tfrac12\rho v^2$: \qty{19.3}{bar} de pressão
e \qty{0.36}{bar} de energia cinética; no bocal a pressão cai
a $P_0$ enquanto $\tfrac12\rho v^2$ sobe a $\rho gH$: logo depois dele, $P = P_0$.

\textbf{8.} $\Delta h = \lambda(\ell/D)v^2/2g = 0.015 \times 133 \times 3.67 = \qty{7.3}{m}$,
$3.7\%$ da potência.

\textbf{9.} $v = \qty{19.1}{m/s}$, $v^2/2g = \qty{18.6}{m}$, $\Delta h = 0.015 \times 200 \times
18.6 = \qty{56}{m}$: $28\%$ perdidos --- a economia de aço é paga a cada
segundo em energia.

\textbf{10.} $v_t = 60/\pi = \qty{19.1}{m/s}$; $\Delta P = 500(19.1^2 - 8.49^2) =
\qty{146}{kPa}$.

\textbf{11.} $\Delta P = (\rho_{\text{Hg}} - \rho)g\,\Delta h$: $\Delta h = 146000/(12600
\times 9.81) = \qty{1.18}{m}$.

\textbf{12.} $\Delta P \propto v^2 \propto D_V^2$: $D_V \propto \sqrt{\Delta h}$; com metade da
vazão dá um quarto, \qty{0.30}{m}.

\textbf{13.} Uma tomada saliente enfrenta ou perturba o escoamento e lê parte
da \omterm{thm:b2:euler-bernoulli:bernoulli}{pressão dinâmica} (um Pitot), e não a pressão estática.

\textbf{14.} $P = P_0 - \rho g \times 6 - \tfrac12\rho v^2 = 100 - 59 - 36 = \qty{5}{kPa}$:
perto da pressão de vapor, o ar dissolvido sai da solução e a água pode
ferver (cavitação), rompendo a coluna. Baixe o ponto alto, ou instale uma
ventosa e uma bomba de vácuo, ou reduza a vazão.

\textbf{15.} $D_V/A = 60/2 \times 10^6 = \qty{3e-5}{m/s} = \qty{11}{cm/h}$.

\textbf{16.} $A\,\dd h/\dd t = -s\sqrt{2gh}$: $\sqrt h = \sqrt{h_0} - (s/A)\sqrt{g/2}\,t$.

\textbf{17.} $\Delta t = (A/s)\sqrt{2/g}(\sqrt{200} - \sqrt{190}) = 2.09 \times 10^6 \times
0.45 \times 0.36 = \qty{3.4e5}{s} = \qty{3.9}{dias}$; na vazão de projeto constante,
$10A/D_V = \qty{3.3e5}{s}$: quase o mesmo, pois o bocal foi dimensionado
para o escoamento de Torricelli com $h \approx H$.

\textbf{18.} $\ell\,\dd v/\dd t = gH - v^2/2$: $v = V\tanh(Vt/2\ell)$, $V = \qty{62.6}{m/s}$,
constante de tempo $2\ell/V = \qty{13}{s}$.

\textbf{19.} Em $v = V/2$: $\dd v/\dd t = (gH - V^2/8)/\ell = \qty{3.7}{m/s^2}$;
$P = P_0 + \rho gH - \tfrac12\rho v^2 - \rho\ell\,\dd v/\dd t = 1.0 + 19.6 - 4.9 - 14.7
= \qty{1.0}{bar}$: a maior parte da carga está ocupada em acelerar a coluna, de modo que
a pressão no fundo fica muito abaixo de seu valor permanente.

\textbf{20.} $|\dd v/\dd t| = 8.49/10 = \qty{0.85}{m/s^2}$: $\Delta P = \rho\ell \times 0.85 =
\qty{3.4}{bar}$.

\textbf{21.} $\Delta P = 1000 \times 1400 \times 8.49 = \qty{119}{bar}$, seis vezes a
pressão estática. A ida e volta da onda leva $2\ell/c = \qty{0.57}{s}$: um
fechamento mais rápido que isso é "súbito" e recebe toda a sobrepressão
de Joukowsky; fechamentos mais lentos recebem menos.

\textbf{22.} $Av = A_s\,\dd z/\dd t$.

\textbf{23.} Euler ao longo do túnel, do reservatório (pressão
$P_0 + \rho g\times$profundidade) até o pé do poço (pressão $P_0 + \rho g
(\text{profundidade} + z)$): $\rho L\,\dd v/\dd t = -\rho gz$; com $v = (A_s/A)\dot z$:
$\ddot z + (gA/LA_s)z = 0$; $T = 2\pi\sqrt{LA_s/gA} = 2\pi\sqrt{2000 \times 100/
68.7} = \qty{340}{s} \approx \qty{5.7}{min}$.

\textbf{24.} $\tfrac12\rho ALv^2 = \tfrac12\rho gA_sz_{\max}^2$ com $v = 60/7 =
\qty{8.6}{m/s}$: $z_{\max} = v\sqrt{AL/gA_s} = \qty{32}{m}$.

\textbf{25.} Barragem: hidrostática, $\rho gh$. Turbina com as válvulas fechadas: o
mesmo, $\rho gH$. Conduto forçado em operação: Bernoulli, $\tfrac12\rho v^2$ tomado
da carga estática. Golpe de aríete: compressibilidade, $\rho c\Delta v$.
Chaminé de equilíbrio: Euler não permanente ao longo do túnel, uma oscilação de período
$2\pi\sqrt{LA_s/gA}$.
\end{solution}
